% from ega4-12-fr.tex
\oldpage[IV-3]{183}
  de type fini, $t$ un élément de l'idéal maximal $\mathfrak m$ de $A$, qui
  n'est contenu dans aucun idéal premier minimal de $A$~; s'il existe un de ces
  idéaux premiers minimaux $\mathfrak p$ tels que $\dim(A/\mathfrak p)=1$,
  est-il vrai que $\mathfrak m$ est un idéal premier associé à $M/tM$ ($t$
  n'étant pas supposé $M$-régulier)~?

  On notera toutefois qu'on ne peut, dans (12.2.1, (xi)), supprimer purement et
  simplement l'hypothèse que $\mathcal F$ est $f$-plat. On prendra en effet ici
  $P$ comme dans la Remarque (ii), et dans $P$ le sous-préschéma fermé $X_0$
  réunion des trois droites $X_1,X_2,X_3$ définies respectivement par
  $t_1=t_3=0$, $t_1-t_0=t_3=0$, $t_2=t_3=0$~; on définit la projection $f_0$ de
  $X_0$ sur $X_1$ comme précédemment et sa restriction $f:X\to Y$, où
  $Y=D_+(t_0)$ est la droite affine et $X=f_0^{-1}(Y)$~; $f$ est propre mais
  non plat, et si l'on prend $\mathcal F=\mathcal O_X$, $\mathcal F_y$ n'a de
  cycles premiers associés immergés pour aucun $y\in Y$. Mais on a $l(y_0)=1$
  et $l(y)=2$ pour $y\neq y_0$ dans $Y$.
\end{enumerate}
\end{remarks}

\begin{theorem}[12.2.4]
\label{IV.12.2.4-fr}
Soient $f:X\to Y$ un morphisme propre, plat et de présentation finie, $k$ un
entier $\geq 1$. Les parties suivantes de $Y$ sont ouvertes~:
\begin{enumerate}[label=(\roman*)]
  \item L'ensemble des $y\in Y$ tels que $X_y$ possède la propriété $(S_k)$.
  \item L'ensemble des $y\in Y$ tels que $X_y$ vérifie la propriété $(R_k)$
  géométrique, soit équidimensionnel en chaque point et n'ait pas de cycle
  premier associé immergé.
  \item L'ensemble des $y\in Y$ tels que $X_y$ soit géométriquement régulier
  (i.e. lisse sur $k(y)$).
  \item L'ensemble des $y\in Y$ tels que $X_y$ soit géométriquement normal.
  \item L'ensemble des $y\in Y$ tels que $X_y$ soit géométriquement réduit.
  \item L'ensemble des $y\in Y$ tels que $X_y$ soit géométriquement réduit et
  que le nombre géométrique de composantes connexes de $X_y$ soit égal à $k$.
  \item L'ensemble des $y\in Y$ tels que $X_y$ soit géométriquement
  ponctuellement intègre (4.6.9).
  \item L'ensemble des $y\in Y$ tels que $X_y$ soit géométriquement intègre.
  \item L'ensemble des $y\in Y$ tels que $X_y$ n'ait pas de cycle premier
  associé immergé et que la multiplicité totale (4.7.4) de $X_y$ sur $k(y)$
  soit $\leq k$.
\end{enumerate}
\end{theorem}

\begin{proof}
Sauf pour (vi), ces assertions sont des cas particuliers d'assertions de
(12.2.1), ou se déduisent d'assertions de (12.1.6) comme au début de la
démonstration de (12.2.1). Pour (vi), on se ramène comme toujours (compte tenu
de l'invariance du nombre géométrique de composantes connexes par changement
de base (4.5.6)) au cas où $Y$ est noethérien. L'ensemble $U$ des $y\in Y$ tels
que $X_y$ soit géométriquement réduit est ouvert dans $Y$ en vertu de (v)~; il
résulte alors de (\textbf{III}, 7.8.7 et 7.8.6) que pour tout $y\in U$, il y a
un voisinage $V\subset U$ de $y$ et un entier $m$ tels que, pour \emph{tout}
$z\in V$, $\Gamma(X_z,\mathcal O_{X_z})$ soit isomorphe à $(k(z))^m$~; mais en
vertu de (\textbf{III}, 4.3.4), $m$ est alors le nombre géométrique de
composantes connexes de $X_z$, d'où la conclusion. On donnera une autre
démonstration de (vi) dans (15.5.9).
\end{proof}

\subsection{Propriétés cohomologiques locales des fibres d'un morphisme plat et
localement de présentation finie}
\label{subsection:IV.12.3-fr}

\begin{lemma}[12.3.1]
\label{IV.12.3.1-fr}
Soient $A$ un anneau, $B$ une $A$-algèbre de présentation finie qui soit un
$A$-module plat, $M$ un $B$-module de présentation finie, qui soit un
$A$-module plat. Alors les conditions suivantes sont équivalentes~:
\begin{enumerate}[label=\emph{\alph*)}]
  \item $M$ est un $B$-module projectif.
  \item $M$ est un $B$-module plat.
  \item Pour tout $y\in\operatorname{Spec}(A)$, $M\otimes_Ak(y)$ est un
  $(B\otimes_Ak(y))$-module projectif.
\end{enumerate}
\end{lemma}

\begin{proof}
